

\section{Introduction} \label{Introduction}
In the previous meeting, it was addressed by the tutor that our physics models are sometimes too complicated or don't match with each other. Therefore, what is the goal of this SSA is based on the models we developed so far, finalise the physics model for the system.

Moreover, because of the confusion caused in my last SSA, not all of the parameters are accurately determined. Therefore, in this SSA, I will finalise the values of those parameters too.

\textbf{[possibly]:} Help Izabella do the simulink simulation.

\section{Model}
\subsection{With or without Thermal energy storage}
 it was mentioned in the previous SSA that without thermal storage, the Isothermal expansion only rely on the environment will not be possible. Because a temperature gradient is needed to let environmental heat flow into the expanding gas. However, the isothermal expansion means the temperature of the gas is constantly equal to the ambient temperature, which means no temperature gradient is created; therefore, no heat flow will occur. 
 
 This statement was not entirely true. Because we could make the process near-isothermal, this means the expanding gas has a slightly lower temperature than the ambient temperature (e.g. a $\Delta T = 20/ 30K$). This will indeed create a temperature gradient, and the heat will indeed flow from the environment into the expanding gas. 
 
\subsubsection{Isothermal vs. Adiabatic}
During expansion phase, the isothermal expansion (which is used by I-CAES) has a better energy extraction rate than adiabatic expansion. Because during adiabatic expansion process, the working gas undergoes expansion without exchanging thermal energy with the surrounding environment ($Q=0$). To perform mechanical work against the turbine blades, the fluid exclusively uses its own internal molecular kinetic (interal) energy, resulting in a rapid decline in fluid temperature, then this thermal reduction causes a accelerated decay in the internal pressure of the gas during expansion. 

In contrast, during an isothermal expansion process, the working air continuously absorbs thermal energy from the thermal energy storage  system. Because the thermal input balances the expansion work, the internal energy of the gas remains constant ($\Delta U = 0$). This sustained temperature prevents the rapid decay of pressure, thereby enabling the fluid to maintain a higher pressure profile throughout the expansion cycle and deliver a greater continuous power output.

The diagram \ref{ada_vs_iso} shows the different between two expansion process.

\subsubsection{Isothermal with TES}
Therefore using isothermal expansion instead of adiabatic makes the system much more efficient. However, the temperature of the expanding gas will affect the output power significantly. For explanation please see the following derivation: 
\begin{figure}[H]
    \centering
    \includegraphics[width=0.5\linewidth]{ada_vs_iso.png}
    \caption{(work is the area under the curve)}
    \label{ada_vs_iso}
\end{figure}

In a closed system, the expansion work performed by a gas is defined by the integral of pressure with respect to volume. For a unit mass of gas ($1\text{ kg}$), this thermodynamic relationship is expressed as
\begin{equation}
w = \int_{v_{in}}^{v_{out}} P  dv.
\end{equation}
**($v$ is specific volume)

According to the ideal gas law, the equation of state for a unit mass is given by $P \cdot v = R_{air} \cdot T$. Consequently, the pressure can be described as $P = \frac{R_{air} \cdot T}{v}$. Substituting this expression into the work integral yields
$$
w = \int_{v_{in}}^{v_{out}} \frac{R_{air} \cdot T}{v} dv.
$$
Since the expansion is modeled as an isothermal process, the temperature $T$ remains constant throughout the state change. The constant terms $R_{air} \cdot T$ can therefore be moved outside the integral, resulting in
$$
w_{iso} = R_{air} \cdot T \int_{v_{in}}^{v_{out}} \frac{1}{v}  dv = R_{air} \cdot T \cdot \ln\left(\frac{v_{out}}{v_{in}}\right).
$$
Finally, Boyle's law dictates that under isothermal conditions, fluid pressure and specific volume are inversely proportional, satisfying $p_{store} \cdot v_{in} = p_{out} \cdot v_{out}$. This yields the volume ratio substitution $\frac{v_{out}}{v_{in}} = \frac{p_{store}}{p_{out}}$. Substituting this relationship into the integrated equation determines the expression for specific isothermal expansion work:
\begin{equation}
w_{iso} = R_{air} \cdot T \cdot \ln\left(\frac{p_{store}}{p_{out}}\right).
\end{equation}
$$P_{discharge} = \dot{m} \cdot w_{iso}$$
So the discharge power equation is:
\begin{equation}
   \fcolorbox{red}{white}{$ P_{discharge} = \dot{m} \cdot R_{air} \cdot T \cdot \ln\left(\frac{p_{in}}{p_{out}}\right)$}
\end{equation}
From \equationautorefname{2.3} we can notice the discharge power is directly proportional to the air temperature. Therefore the presence of thermal energy storage is important to maintain the necessary efficiency of our system. 

**Actually in reallity achieving a perfectly isothermal process is practically impossible due to the limited residence time of the air and the finite heat transfer rate. In reality, the expansion is a near-isothermal polytropic process with a polytropic index $n$ close to 1 (e.g., $n \approx 1.05 \sim 1.1$). However, the fundamental thermodynamic principle remains. The actual model we use in simulation will be the polytropic one, so we have more accurate model, and we can continue to use the experiment plan Seb and Alex developed.
\subsubsection{Isothermal without TES}
BUT do we really need TES to boost up the discharging power?

First we can analyse it \textbf{from the discharge power equation aspect}. 

The maximum discharge power of the demand is around 70 MW, and the storage pressure is 80 bar. So the discharging equation would be like:
$$70,000,000 = \dot{m} \cdot 287 \cdot T \cdot \ln\left(\frac{80}{1}\right)$$
From where we get $$\dot{m} = \frac{55,661}{T}$$
From changing the value of the T, we could notice if we keep the expanding air under temperature $T = 250 \text{ K}$, The mass flow rate would be $\approx 222.64 \text{ kg/s}$, And if we heat the air up to 500K, the needed mass flow rate would be $\approx 113 \text{ kg/s}$

The diagram \ref{fig:Mass flow rate versus temperature} shows the relation of two parameters. 
\begin{figure}[H]
    \centering
    \includegraphics[width=0.5\linewidth]{Mass flow rate versus temperature.png}
    \caption{Mass flow rate versus temperature}
    \label{fig:Mass flow rate versus temperature}
\end{figure}
From the sources both 223 and 113 are reasonable numbers for a CAES (e.g. Huntorf CAES Mass Flow Rate aroung 417 kg/s). So from the discharge power equation aspect, the temperature does not change the mass flow rate significantly. 
\vspace{1em}

Then we look at what \textbf{physics phenomenon are governing the discharging} when the expanding air is heated up by the environment only.

The discharging power should be equal to the heat exchange, so the value of heat exchange must keep up with the demand. 
\begin{equation}
P_{discharge}\approx\dot{Q}_{amb} = U \cdot A \cdot (T_{amb} - T_g) = U \cdot A \cdot \Delta T
\end{equation}
where $\dot{Q}_{amb}$ represents the thermal power transferred from the environment to the expanding air, $U$ is the overall heat transfer coefficient, $A$ is the heat transfer surface area, $T_{amb}$ is the ambient temperature, $T_g$ is the localized temperature of the expanding gas, and $\Delta T$ represents the temperature difference.

If we allow a Delta T of 30K, the value of $U\cdot A$ would be 2 333 333, the relation between U and A would be $U=\frac{2333333}{A}$

Typical $U$ Value for Air-to-Air Heat Exchangers is about $20 - 50 \text{ W/m}^2\text{K}$, which means around $77,777 m^2$ of surface area is needed. Which is seemingly impossible to achieve.

Another alternative is to design our custom  heat exchanger to boost U value up and increase surface area, but I think this is complex because the device needs to conduct air expansion and heat exchange simultaneously while having a low working power. And also I am not sure if it's allowed to do in this project.  Also the temperature in the winter is much lower than in summer so it will decrease the delta T even more. So the limit of this method is not with the mass flow rate, (although the increased mass flow rate does mean the discharging is less efficient), it is with the heat exchange rate. That's why \textbf{in summary the isothermal compression with TES is the most feasible solution for our storage system.} 

\vspace{1em}

\textbf{An idea for experiment team:} maybe a good validation to do is prove in real life the difference of the system with TES and without TES. 

\subsection{Listing the equations of model}
 so now we have proven that the isothermal compression with the thermal energy storage is the most feasible solution, we could continue working on the models. In this part I will summarise the equations we have determined from each team, evaluate them, unify them, or maybe modify them so we can have one mathematical model for the system. And they are listed below:
 
 \subsubsection{Transporting from grid}
 \begin{equation}
P_{net} = P_{supply} - P_{demand}
\end{equation}
\vspace{1em}

\subsubsection{Injection \& Compression}
\begin{equation}
P_{eff}(t) =\eta_{tran} \cdot
\begin{cases}
\min(P_{net}, P_{limit}) & \text{if } P_{net}(t) \ge 0
\cr P_{net}(t) & \text{if } P_{net}(t) < 0
\end{cases}
\end{equation}
\vspace{1em}

\begin{equation}
E_{eff} = P_{eff}, dt
\end{equation}
 %\int \eta_{in} \cdot
 Through compressors, the electricity energy will be stored in two means: Thermal energy and Pressure Exergy (about exergy, see: \href{https://en.wikipedia.org/wiki/Exergy}{wikipedia})

 The polytropic compression work:
$$ P_{charge} =\eta_{tran} \cdot
\begin{cases}
\min(P_{net}, P_{limit}) & \text{if } P_{net}(t) \ge 0
\cr 0 & \text{if } P_{net}(t) < 0
\end{cases}$$

\begin{equation}
P_{charge} = \eta_{comp} \cdot\frac{n}{n-1} \dot{m}_{charge} R_{air} T_{in} \left[ \left( \frac{p_{store}}{p_{amb}} \right)^{\frac{n-1}{n}} - 1 \right]
\end{equation}

By rearranging the equation, we can dynamically calculate the air mass flow rate that the system can inject into the underground gas storage chamber at any given time $t$:

\begin{equation}
    \dot{m}_{charge} = \frac{P_{charge}}{ \eta_{comp} \cdot\frac{n}{n-1} R_{air} T_{in} \left[ \left( \frac{p_{store}}{p_{amb}} \right)^{\frac{n-1}{n}} - 1 \right]}
\end{equation}
Where:
\begin{itemize}
    \item[] $n = \text{Polytropic index (-)}$ (for near-isothermal process, $n \approx1.05-1.2$
    \item[] $\dot{m}_{charge} = \text{Mass flow rate of charging (kg/s)}$
    \item[] $R_{air} = \text{Specific gas constant of air (J/(kg} \cdot \text{K))}$
    \item[] $T_{in} = \text{Ambient inlet temperature (K)}$
    \item[] $p_{amb} = \text{Atmospheric pressure (Pa)}$
    \item[] $p_{store} = \text{Target storage pressure in the cavern (Pa)}$
\end{itemize}



 \textbf{Thermal power} (derivation in \ref{append: Derivation of Thermal power}):
 \begin{equation}
     P_{thermal}(t) = \eta_{thermal} \cdot [P_{charge}(t) - \dot{m}_{charge}(t) \cdot c_p \cdot (T_{out} - T_{in})]
 \end{equation}
 
 - thermal efficiency factor ($\eta_{thermal}$) is the fraction of the input compression work that becomes usable thermal energy :

$$ \eta_{thermal} = \frac{P_{thermal}}{P_{charge}} = 1 - \frac{c_p (T_2 - T_1)}{w_{comp}} $$
$$ w_{comp} = \frac{P_{charge}}{\dot{m}_{charge}}  $$

- $T_{out}$ is the actual air temperature at the compressor outlet, which can be calculated using the polytropic relation
$$T_{out} = T_{in} \left(\frac{p_{store}}{p_{amb}}\right)^{\frac{n-1}{n}}$$
\textbf{Pressure Exergy Rate} (derivation in \ref{append:Derivation of pressure exergy rate}):
\begin{equation}
    P_{exergy}(t) = \dot{m}_{charge}(t) \cdot R_{air} \cdot T_{in} \cdot \ln\left(\frac{p_{store}(t)}{p_{amb}}\right)
\end{equation}

\subsubsection{Storage systems - thermal decay}
 as mentioned in the previous SSA the pressure decay of CAES is extremely low and negligible.  Therefore we only need a model for thermal decay in the TES.
The heat stored in the surface TES will inevitably be lost to the environment over time. This thermal decay over time is primarily governed by Newton's Law of Cooling.

Instantaneous Heat Loss Rate:
\begin{equation}
\dot{Q}_{loss}(t)=U_{tes}A_{tes}(T_{tes}(t)-T_{amb})
\end{equation}
where
\begin{itemize}
    \item []$U_{tes}$ is the overall heat transfer coefficient of the thermal energy storage tank
    \item[] $A_{tes}$ is the surface area
    \item[] $T_{amb}$ is the ambient temperature.
\end{itemize}  

Assuming the medium inside the TES is well-mixed, its temperature change rate is given by:
\begin{equation}
m_{tes}c_{tes}\frac{dT_{tes}}{dt}=-U_{tes}A_{tes}(T_{tes}(t)-T_{amb})
\end{equation}
where 
\begin{itemize}
    \item []$m_{tes}$ is the mass of the storage medium
    \item [] $c_{tes}$ is the specific heat capacity
\end{itemize}

By integrating the above equation, the formula for the exponential decay of the storage medium's temperature over time is obtained:
\begin{equation}
T_{tes}(t)=T_{amb}+(T_{initial}-T_{amb})e^{-\frac{U_{tes}A_{tes}}{m_{tes}c_{tes}}t}
\end{equation}

\textbf{Storage volume: }

\begin{equation}
    E(t)=\int P_{charge} dt
\end{equation}

\subsubsection{Extraction stage}
When high-pressure air expands to perform work in a turbine, the air temperature drops sharply if there is no external heat supplementation (it will then be adiabatic process). To maintain high power output, the system directly injects the hot fluid stored in the TES into the expansion chamber.

Similar to the charging equation, this process will also be polytropic:
\begin{equation}
    P_{discharge}(t) = \eta_{exp} \cdot \frac{n}{n-1} \cdot \dot{m} \cdot R_{air} \cdot T_{in} \left[ 1 - \left( \frac{p_{amb}}{p_{in}} \right)^{\frac{n-1}{n}} \right]
\end{equation}

$$T_{out} = T_{in} \left( \frac{p_{amb}}{p_{store}(t)} \right)^{\frac{n-1}{n}}$$

Ideally, the $T_{out}=T_{amb}$ to have maximum effciency.

\vspace{1em}
\textbf{Determine n}

We already know the temperature-pressure relation for polytropic expansion:
\begin{equation}
\frac{T_{out}}{T_{in}} = \left( \frac{p_{amb}}{p_{store}(t)} \right)^{\frac{n-1}{n}}
\end{equation}
To solve for $n$, we take the natural logarithm of both sides:
\begin{equation}
\ln\left(\frac{T_{out}}{T_{in}}\right) = \frac{n-1}{n} \cdot \ln\left(\frac{p_{amb}}{p_{store}}\right)
\end{equation}
Through rearrangement, we can isolate $n$ as an independent variable:
\begin{equation}
n = \frac{\ln\left(\frac{p_{amb}}{p_{store}}\right)}{\ln\left(\frac{p_{amb}}{p_{store}}\right) - \ln\left(\frac{T_{out}}{T_{in}}\right)}
\end{equation}

\vspace{1em}

To maintain this polytropic exponent $n$, the thermal energy storage system (TES) must inject thermal power $\dot{Q}_{in}$ into the turbine:
\begin{equation}
\dot{Q}_{in}(t) = \frac{P_{discharge}(t)}{\eta_{exp}} - \dot{m} \cdot c_p \cdot (T_{in} - T_{out})
\end{equation}

\subsubsection{Transporting to demand}
\begin{equation}
    P_{output}=\eta_{tran} \cdot P_{discharge}
\end{equation}

\subsubsection{Other equations}
\begin{equation}
\dot{m} = \rho v A
\end{equation}

\begin{equation}
\rho = \frac{p_{\text{in}}}{R T_{\text{in}}}
\end{equation}

\begin{equation}
A = \frac{\pi}{4} D^2
\end{equation}
 

\subsection{Energetics and exergetics analysis of the system}
 this part of the SSA may help to understand the system from a different perspective. 

 \vspace{1em}
In a thermal energy storage framework, system are often simply described by conservation of energy. However, the I-CAES operates on a more complex theory. In this system the electrical energy is first divided into exergy and energy during the compression stage then reunified during the extraction stage. Compared to the traditional thermal energy storage system, this system circumvents the Carnot efficiency limits. \href{https://news.mit.edu/2010/explained-carnot-0519}{[Explained: The Carnot Limit]} \href{https://en.wikipedia.org/wiki/Carnot%27s_theorem_(thermodynamics)0}{[Carnot's theorem]}

The working mechanism is comprehensively described with the first and second law of thermodynamics.

During the compression the energy of electricity is converted into two different forms: the Microscopic kinetic energy our so-called thermal energy is going to be stored in TES, the thermal energy storage. The compressed high-pressure air will contain a so-called exergy \href{https://www.youtube.com/watch?v=2TPefBIVZZo&pp=ygUGZXhlcmd5}{[Exergy Vs Energy]}. This is the mechanical potential of the system, and it is stored in the extreme disequilibrium of pressure. 

If we expand the high-pressure air without supplementing thermal energy to it, the process will be adiabatic. And the energy we can extract from it will be low. Because we are basically extracting the energy from the temperature remaining in the air, which is just room temperature. More importantly the expanding air will be so cold, that freezes the blades of machine. 

In most cases exergy and energy are stored in the same form in the same medium, for example spring. They are both stored in the deformation. But in our system they are decoupled and stored separately, the energy is in the thermal tank and the exergy is in the pressure.  

When we expand the air, we kind of couple these two back together, the thermal energy is injected back to the air, and the air transfers this energy to the generator while expanding. The benefit of this, as mentioned before, has a significantly better energy efficiency compared to heat engine, as in this case it is not doing the Carnot cycle, therefore not limited by Carnot efficiency limit. 

\vspace{2em}

Another way to understand this would be:
 if we make the system adiabatic,  
 when compressing the air, we are using the electrical energy to do work to compress the air. If we do not collect the heat, it will go to the environment, and the temperature of air will eventually be room temperature in a high pressure. This means we lost the energy but we stored the exergy, during the extraction phase, the high-pressure air still has the tendency to expand on its own, but now it must use it's remaining internal energy (which is stored in its room temperature). This means it will cool down significantly, in theory because the internal energy of an ideal gas is directly proportional to the temperature, it will cool down by the same amount of temperature as it heats up. 


\section{Finalizing Parameters}
 in this part summirizes and lists the parameters needed for modelling, some of them are already determined before. Some of them are being determined now.

 \vspace{1em}

\subsection{Constants}
Ambient temperature: $T_{amb}= 300 \text{ (K)}$  
Polytropic index 

Specific gas constant of air: $R_{air}= 287.05 \text{ (J/(kg} \cdot \text{K)})$

 Specific heat capacity of air at constant pressure: $c_p= 1005 \text{ (J/(kg} \cdot \text{K)})$ 

Specific heat capacity of water (The thermal tank medium): $c_{tes}= 4184 \text{ (J/(kg} \cdot \text{K)})$

\subsection{ System Design Parameters}
Maximum allowable power draw from the grid for charging: $P_{limit}=300(\text{MW})$

Target maximum storage pressure in the underground cavern: $p_{store}=80\cdot10^5\text({Pa})$

Overall heat transfer coefficient of the thermal energy storage tank: $U_{tes}=(\text{W/(m}^2 \cdot \text{K)})$

Surface area of the thermal energy storage tank: $A_{tes}=(\text{m}^2)$

Total mass of the storage medium within the TES: $m_{tes}=(\text{kg})$

Diameter of the injection/extraction piping: $D=(\text{m})$

Cross-sectional area of the piping: $A=(\text{m}^2)$

Maximum allowable mass flow rate: $\dot{m}_{max}=$

TES internaland compression temerature : $T_{in}$

Temperature we are heating the expanding air: $T_{heat}$

\subsubsection{Storage Volume}
\textbf{Energy needs to be stored:}

\textbf{Cavern total volume}

\textbf{TES total volume:}

\subsection{Efficiency coefficients}

\begin{itemize}
    \item Transmission efficiency (accounts for losses when transporting power from/to the grid): $\eta_{tran}=0.97$
    \item Compressor mechanical/isentropic efficiency: $\eta_{comp}=0.833$
    \item Thermal efficiency factor (fraction of compression work successfully converted into usable thermal energy): $\eta_{thermal}=(-)$
    \item Expansion turbine efficiency: $\eta_{exp}=(-)$
\end{itemize}


 

  
\section{MatLab implementation}
%4. 状态变量积分 (Integration of State Variables)为了在 Simulink 等动态仿真软件中追踪系统的整体储能状态，我们需要对上述瞬时功率随时间进行积分（integrals）。累计储存的热能 (Total Stored Thermal Energy in TES):$$E_{thermal}(t) = E_{thermal}(0) + \int_{0}^{t} P_{thermal}(\tau) \, d\tau$$累计储存的压力㶲 (Total Stored Pressure Exergy in Cavern):$$E_{exergy}(t) = E_{exergy}(0) + \int_{0}^{t} P_{exergy}(\tau) \, d\tau$$同时，储气室内部的总质量（Total Air Mass）也在累积，这将直接用于更新下一时刻的 $p_{store}(t)$：$$M_{cavern}(t) = M_{cavern}(0) + \int_{0}^{t} \dot{m}_{charge}(\tau) \, d\tau$$
\section{Conclusion}